% 00_checklist.tex — чек-лист проверки документов команды перед сдачей.
% Собрать: ./build.sh docs/00_checklist.tex. Автоматическая часть раздела «Идентификация»: ./build.sh --check
\documentclass{tpks}
\doctitle{Чек-лист проверки документов}
\docversion{1}
\docdate{20.11.2026}

\begin{document}
\makeheader
\begin{changelog}
\change{1}{20.11.2026}{Первая версия: общие проверки, идентификация требований, требования, процессы, управление конфигурацией, оформление}{Чернов А. В.}{Замечания преподавателя}
\end{changelog}

\section{Введение}

\subsection{Назначение}
Чек-лист перечисляет, что проверить в документе перед каждой сдачей. Он собран из чек-листа, по которому команда готовилась к сдачам, и из замечаний преподавателя к нашим документам. Главный раздел --- проверка идентификации требований (раздел~\ref{sec:ids}).

\subsection{Область применения}
Чек-лист применяют ко всем документам команды. Пункты, которые к документу не относятся, отмечают «не применимо»: например, раздел~\ref{sec:proc} не относится к спецификациям требований.

\subsection{Порядок проверки}
\begin{enumerate}
  \item Автор собирает документ командой \texttt{./build.sh} и запускает \texttt{./build.sh -{}-check}. Ошибки исправляют до ручной проверки.
  \item Проверяющий, который не писал документ, проходит все пункты и отмечает: «да», «нет» или «не применимо».
  \item Каждое «нет» исправляют или записывают как открытый вопрос.
  \item Для повторной сдачи собирают экземпляр с выделенными изменениями: \texttt{./build.sh -{}-changes}.
\end{enumerate}

В столбце «Как» таблиц указано: «сборка» --- проверку выполняет шаблон при сборке PDF; «скрипт» --- \texttt{./build.sh -{}-check}; «вручную» --- проверяющий.

\section{Общие проверки}
\begin{tabl}{|c|L|P{2.2cm}|}{Общие проверки}[tab:common]{\textbf{№} & \textbf{Проверка} & \textbf{Как}}
О1 & Термины, процессы, входы и выходы одинаковы во всех документах команды. Одно понятие называется везде одинаково & вручную \\ \hline
О2 & Определён каждый термин, который можно понять неоднозначно, каждое сокращение и каждая единица измерения & вручную \\ \hline
О3 & Термин «формальная инспекция» определён в глоссарии: в авиационном стандарте его нет & вручную \\ \hline
О4 & Во введении есть «Назначение», «Область применения», «Термины и сокращения», «Ссылки» & вручную \\ \hline
О5 & Термины Git (main, коммит, pull request) встречаются только в Плане управления конфигурацией ПО; в остальных документах --- «базовая версия», «ЗИ», «СП» & вручную \\ \hline
О6 & Нет ссылок на слайды, конспект лекций, гайд старших курсов, рабочие материалы и пути к файлам (кроме мест хранения в Плане УК ПО) & вручную \\ \hline
О7 & Нет противоречий между разделами документа и с другими документами команды & вручную \\ \hline
\end{tabl}

\section{Идентификация требований}\label{sec:ids}
Раздел относится к документам с требованиями и к документам, которые на требования ссылаются. Идентификатор записывают как «ПРЕФИКС\_ФФ\_НН; \#версия»: префикс SYS или SRD, номер функции, порядковый номер, служебные символы «;~\#» и номер версии требования.

\begin{tabl}{|c|L|P{2.2cm}|}{Проверки идентификации}[tab:idcheck]{\textbf{№} & \textbf{Проверка} & \textbf{Как}}
И1 & Каждый идентификатор записан по схеме SYS\_ФФ\_НН или SRD\_ФФ\_НН: две цифры функции и две цифры номера & сборка, скрипт \\ \hline
И2 & Идентификатор уникален. Если требование приведено в другом документе как образец, у него тот же текст и та же версия & сборка, скрипт \\ \hline
И3 & Версия требования --- целое число от 1. Заголовок требования печатается со служебными символами «ID; \#версия», ссылка на родительское требование в поле «Источник» --- тоже «ID; \#версия» & сборка, скрипт \\ \hline
И4 & Требования объявлены в документе по возрастанию идентификаторов: SRD\_11 не стоит раньше SRD\_01 & скрипт \\ \hline
И5 & Внутри функции номера идут с 01 без пропусков. Пропуск допустим, только если требование удалено & скрипт \\ \hline
И6 & По сравнению с принятой версией документа: ни один идентификатор не пропал без записи в листе регистрации; у каждого требования с изменённым текстом версия увеличена; версия нигде не уменьшилась & скрипт \\ \hline
И7 & Номер функции требования к ПО совпадает с номером функции хотя бы одного его системного требования & скрипт \\ \hline
И8 & У каждого требования к ПО указан источник: родительские системные требования или отметка «производное». Каждый родитель существует в «Системных требованиях» & скрипт \\ \hline
И9 & Каждое системное требование покрыто хотя бы одним требованием к ПО & скрипт \\ \hline
И10 & Матрица трассировки совпадает с полями «Источник»: те же пары в обе стороны & скрипт \\ \hline
И11 & Производные требования помечены, у каждого в комментарии есть обоснование & скрипт, вручную \\ \hline
И12 & Сводная таблица идентификаторов совпадает с документом: диапазоны существуют, количества верны, каждое требование попало в одну строку & скрипт \\ \hline
И13 & Каждая ссылка на требование в тексте и таблицах указывает на существующий идентификатор & скрипт \\ \hline
И14 & В документе расшифрованы префиксы и части идентификатора (SYS, SRD, ФФ, НН, версия), есть сводная таблица идентификаторов & вручную \\ \hline
И15 & Идентификатор удалённого требования не использован для нового; удаление записано в листе регистрации & вручную \\ \hline
\end{tabl}

\section{Требования}
\begin{tabl}{|c|L|P{2.2cm}|}{Проверки требований}[tab:reqcheck]{\textbf{№} & \textbf{Проверка} & \textbf{Как}}
Т1 & В каждом требовании есть «должен», «должна» или «должно» & вручную \\ \hline
Т2 & Требование говорит, ЧТО делает система или ПО, а не КАК: нет технологий, кнопок, окон и деталей реализации & вручную \\ \hline
Т3 & Требование проверяемо: указан способ проверки и критерий «пройдено / не пройдено» & вручную \\ \hline
Т4 & Требование сформулировано положительно: «должна не …», а не «не должна» & вручную \\ \hline
Т5 & Требование неделимо и однозначно, каждое встречается в документе один раз & вручную \\ \hline
Т6 & В требованиях высокого уровня нет компонентов ПО и деталей реализации, кроме проектных ограничений & вручную \\ \hline
Т7 & Нет TBD без причины и условия, при котором TBD снимается & вручную \\ \hline
Т8 & Комментарии отделены от текста требования и оформлены в поле «Комментарий» & вручную \\ \hline
Т9 & Числовые характеристики заданы числами с единицами, у входов указаны допустимые диапазоны & вручную \\ \hline
\end{tabl}

\section{Процессы}\label{sec:proc}
Раздел относится к планам и стандартам.

\begin{tabl}{|c|L|P{2.2cm}|}{Проверки описания процессов}[tab:proccheck]{\textbf{№} & \textbf{Проверка} & \textbf{Как}}
П1 & Описаны девять процессов жизненного цикла, интеграция --- отдельный процесс & вручную \\ \hline
П2 & У каждого процесса есть цели, входные данные, условия входа, описание, условия выхода, выходные данные и роли & вручную \\ \hline
П3 & Выход кодирования --- исходный код и СП; исполняемый объектный код --- выход интеграции & вручную \\ \hline
П4 & У процессов разработки ЗИ среди входов, СП среди выходов; план и соответствующий стандарт среди входов & вручную \\ \hline
П5 & Входы разработки требований к ПО: системные требования, архитектура системы, программно-аппаратные интерфейсы. Выходы: спецификация требований к ПО и СП & вручную \\ \hline
П6 & Входы интеграции: исходный код, параметры сборки, описание проекта. Выходы: исполняемый объектный код, протокол сборки, командные файлы, протокол загрузки & вручную \\ \hline
П7 & Для каждого результата указано конкретное мероприятие верификации: рассмотрение, анализ или тестирование; для спецификации требований --- формальная инспекция & вручную \\ \hline
П8 & Нигде не написано о верификации процессов УК, обеспечения качества и взаимодействия с сертифицирующим органом & вручную \\ \hline
П9 & Проверка системных требований называется валидацией, требований к ПО --- верификацией & вручную \\ \hline
П10 & Соблюдена независимость: проверяет не автор, автор кода не пишет тесты, разработчики не инспектируют свой результат & вручную \\ \hline
\end{tabl}

\section{Управление конфигурацией}
\begin{tabl}{|c|L|P{2.2cm}|}{Проверки управления конфигурацией}[tab:scmcheck]{\textbf{№} & \textbf{Проверка} & \textbf{Как}}
К1 & Категории контроля КК1 и КК2 присвоены конкретным объектам, а не проекту целиком & вручную \\ \hline
К2 & Объект КК1 после базовой версии меняется только по ЗИ & вручную \\ \hline
К3 & Базовая версия появляется после формальной инспекции без открытых несоответствий & вручную \\ \hline
К4 & Описан путь «ошибка → СП → совет по изменениям → ЗИ → изменение → проверка → новая базовая версия» & вручную \\ \hline
К5 & Описано, где хранится каждый результат и кто кому его передаёт & вручную \\ \hline
\end{tabl}

\section{Оформление}
Оформление проверяют по таблице «Что проверяет автор» из «Требований к оформлению документов». Дополнительно:

\begin{tabl}{|c|L|P{2.2cm}|}{Проверки оформления}[tab:formcheck]{\textbf{№} & \textbf{Проверка} & \textbf{Как}}
Ф1 & Порядок частей: титульный лист, лист регистрации изменений, содержание, основной текст & сборка \\ \hline
Ф2 & Версия и дата на титульном листе совпадают с последней строкой листа регистрации & вручную \\ \hline
Ф3 & На каждую таблицу и каждый рисунок есть ссылка в тексте, все подписаны & вручную \\ \hline
Ф4 & Для повторной сдачи собран экземпляр с выделенными изменениями & вручную \\ \hline
\end{tabl}

\end{document}
